\scanpage{176}\mspage{31}%
haut s'interprète immédiatement comme
\[
\varphi(F_1)\simeq p_*j^*(\widetilde F_1)
\]
où $F_1\mapsto\widetilde F_1$ est la restriction de $X_1$ à $\widetilde X_1$ (qui est donc une équivalence de catégories pour les faisceaux (loc. const.) envisagés). \deleted{\unsure{Il|Il s'|j} \unsure{peut|peut se|\ensuremath{j^*}}} \inserted{La donnée} d'un\inserted{\unsure{e}} $F=(F_0,F_1,u:F_0\to\varphi F_1)$ peut donc aussi s'expliciter comme la donnée de
\[
F_0,\ \widetilde F_1,\ \widetilde u:p^*(F_0)\longrightarrow j^*(\widetilde F_1)
\]
où $F_0$ ($\widetilde F_1$) est un faisceau loc. constant sur $X_0$ (resp. $\widetilde X_1$). Il faut se rappeler de plus qu'ici $p$ est une vraie fibration, et $j$ une inclusion (et pratiquement, -- cas $X_0,X_1$ lisses -- \deleted{\unsure{une}} l'inclusion du bord dans une variété à bord).

Si on préfère, on peut aussi prendre le diagramme un peu moins joli (mais un peu plus canonique)
\diagram{\begin{tikzcd}[ampersand replacement=\&,column sep=large,row sep=large]
\mathring T \arrow[r,"\unsure{J|j}'"] \arrow[d,"p'"] \& X_1\\
X_0\raisebox{0pt}[0pt][0pt]{\parbox{0pt}{\figure{\begin{tikzpicture}[overlay]
% Isolated dark short stroke, whole-image approximately (352,1847), left of X_0; origin uncertain.
\draw[line width=.45pt,line cap=round] (-3.4cm,.08cm) -- (-3.38cm,.03cm) -- (-3.38cm,-.02cm);
\end{tikzpicture}}}} \&
\end{tikzcd}}
revenant essentiellement au $\unsure{\overline m|=}$, car formé d'espaces homotopes au précédent. On peut $\overline m$ y remplacer $X_0$ par $T$ ($X_0$ en étant un rétract par déformation) et écrire
\diagram{\begin{tikzcd}[ampersand replacement=\&,column sep=large,row sep=large]
\mathring T \arrow[r,hook,"J'"] \arrow[d,hook,"p''"] \& X_1\\
\unsure{T|I} \&
\end{tikzcd}}
% The entire upper row and both arrows of the second diagram are visible on this page; its lower target is crowded against the paper edge.
% archive sheet number: 176
